% =====================================================================
%  Resumen
% =====================================================================
\chapter*{RESUMEN}

Se propone el desarrollo de un sistema de comunicación asistida mediante señales
electroencefalográficas (EEG), empleando el paradigma P300 Speller con una matriz de
caracteres que se ilumina sistemáticamente para la clasificación del potencial de eventos
P300 mediante modelos de aprendizaje profundo, permitiendo la identificación del carácter
objetivo. Está dirigido a personas con compromisos motores severos. Así mismo, se
evaluarán estrategias para optimizar la practicidad de la comunicación mediante la
reducción del número de canales y repeticiones requeridas, sin comprometer el desempeño
del modelo, con el propósito de mejorar la eficiencia en la comunicación de los usuarios.

\vspace{0.3cm}
\noindent\textbf{Palabras clave -} Aprendizaje profundo, Arquitecturas neuronales,
Procesamiento de señales, Reconocimiento de patrones.
